% Student assessment reconstructed from the complete matching solution PDF.
% Build twice for question references; see the build guide for TEXINPUTS.
\makeatletter
\def\input@path{{preamble/}{../preamble/}{../../preamble/}{../../../preamble/}{../../../../book2026-2027/preamble/}}
\makeatother
\input{preamble_exams}
\graphicspath{{./}{preamble/}{../preamble/}{../../preamble/}{../../../preamble/}{../../../../book2026-2027/preamble/}}

% Assessment-specific settings.
\setYColumnWidth{0.42}
\newcommand{\CriterionA}{C1}
\newcommand{\CriterionB}{C2}
% Ledger: C1 = 4+4 = 8; C2 = 2+2+4+6 = 14; overall = 22.

% Preserve starred/unstarred calls and subtract color-box padding ONCE.
% Do not reduce linewidth again at individual calls.
\let\ExamSubsectionBox\SubsectionBox
\RenewDocumentCommand{\SubsectionBox}{s m}{%
  \par% Column breaks are placed between complete questions below.
  \begingroup
  \setlength{\linewidth}{\dimexpr\linewidth-2\fboxsep\relax}%
  \IfBooleanTF{#1}{\ExamSubsectionBox*{#2}}{\ExamSubsectionBox{#2}}%
  \endgroup\nobreak
}
% Shared headings use stepcounter, not refstepcounter: capture the automatic
% number explicitly for score-table references, without manual question IDs.
\makeatletter
\newcommand{\QuestionLabel}[1]{%
  \begingroup\edef\@currentlabel{\arabic{subsectionboxproblem}}%
  \label{#1}\endgroup
}
\makeatother
\linespread{1.0}

\setlength{\emergencystretch}{1.5em}
\begin{document}
\noindent
\begin{tabularx}{\textwidth}{|C{2.8cm}|C{\dimexpr\textwidth-5.6cm-6\tabcolsep-4\arrayrulewidth\relax}|C{2.8cm}|}
\hline
\centering\vspace{3mm}\includegraphics[width=2.5cm]{logo.png}
& \CellCenter{\textbf{GLOBAL ECONOMICS}\par
\textbf{Grade 10 -- Learning Evidence 1}\par
\textbf{Discrete monetary outcomes: PMFs and CDFs}\par
\textbf{Teacher: Nicolás López Cuéllar}}
& \CellCenter{\textbf{TERM 2}\par\textbf{2026--2027}}\\\hline
\end{tabularx}\par\smallskip
\LearningObjective{Model discrete monetary outcomes in economic and financial contexts; derive and verify probability mass functions (PMFs), and construct and interpret cumulative distribution functions (CDFs) with tables and graphs.}{%
\textbf{C1}: Characterises discrete random variables in economic and financial contexts by identifying their support, probability mass function and possible outcomes.\par
\textbf{C2}: Interprets PMFs and CDFs by verifying their defining properties and using them to calculate point, interval, cumulative, and complementary probabilities in financial contexts.
}\par\smallskip
\StudentInfo\par

\begin{multicols*}{2}
\raggedcolumns
\setlength{\parskip}{3pt}
\SubsectionBox*{Criteria assessment}

This exam assesses C1 and C2. A \textbf{minimum of 4 points per criterion} is required to pass. Show all relevant calculations and justify your answers. Partial credit may be awarded for correct procedures despite arithmetic errors; answers without demonstrated procedures receive no credit. A 1-point penalty per criterion applies if the student's name is missing or the assessment is submitted late.

\noindent\begin{minipage}{\linewidth}
\centering\small
\setlength{\tabcolsep}{6pt}\renewcommand{\arraystretch}{1.05}
\textbf{Available points}\par\smallskip
\begin{tabular}{|c|c|c|}
\hline\textbf{\textnumero} & \textbf{C1} & \textbf{C2} \\ \hline
\ref*{q:1} & 4 & --- \\\hline
\ref*{q:2} & 4 & --- \\\hline
\ref*{q:3} & --- & 2 \\\hline
\ref*{q:4} & --- & 2 \\\hline
\ref*{q:5} & --- & 4 \\\hline
\ref*{q:6} & --- & 6 \\\hline
\textbf{Total} & \textbf{8} & \textbf{14} \\\hline
\end{tabular}
\end{minipage}\par
\SubsectionBox{Online order contribution\score{C1-\sml{4}}}
\QuestionLabel{q:1}
A small seller handles three orders. Each order has exactly two mutually exclusive possible outcomes: the buyer keeps it with probability 0.6, or cancels it with probability 0.4. These probabilities apply to every order. Each kept order contributes 5 USD and each cancellation contributes no money. A fixed handling cost of 4 USD is then deducted once from the combined contribution. There are no other receipts or costs. Record the net contribution for the three orders.

\begin{enumerate}[label=(\alph*)]
\item List every complete outcome in order and calculate its probability. \sscore{\sml{1}}
\item Define the random variable and its units. Map every complete outcome to a value, identify the support, and construct the grouped PMF. \sscore{\sml{2}}
\item Calculate and interpret the probability that the net contribution is at least 6 USD. \sscore{\sml{1}}
\end{enumerate}

\columnbreak
\SubsectionBox{Two supplier loyalty credits\score{C1-\sml{4}}}
\QuestionLabel{q:2}
A retailer receives one loyalty credit from each of two suppliers. The suppliers award credits independently. Each has exactly three mutually exclusive possible outcomes: no credit, a 3 USD credit, or an 8 USD credit. For the first supplier these probabilities are 0.2, 0.5 and 0.3, respectively. For the second supplier they are 0.4, 0.4 and 0.2, respectively. Add both credits; there are no other adjustments. Record the total credit, keeping first-supplier then second-supplier order.

\begin{enumerate}[label=(\alph*)]
\item List every complete ordered outcome and calculate its probability. \sscore{\sml{1}}
\item Define the random variable and units; map every outcome to a value; identify the support; construct the grouped PMF. \sscore{\sml{2}}
\item Calculate and interpret the probability that the total credit is at least 11 USD. \sscore{\sml{1}}
\end{enumerate}

\SubsectionBox{Local parcel delivery fees\score{C2-\sml{2}}}
\QuestionLabel{q:3}
A cooperative sends one parcel using exactly one of eight exhaustive, mutually exclusive delivery methods: walking courier, bicycle, cargo bicycle, motorcycle, car, van, refrigerated van or truck. Their probabilities are 0.10, 0.10, 0.15, 0.20, 0.15, 0.10, 0.10 and 0.10, respectively. The corresponding fixed handling fees are 0, 0, 2, 4, 7, 11, 16 and 23 USD, respectively. No other handling charges apply. Record the handling fee for the parcel.

\begin{enumerate}[label=(\alph*)]
\item Define the random variable and support; derive the grouped PMF; produce its table and a labeled graph with actual monetary spacing. \sscore{\sml{1}}
\item Construct the complete CDF table, including both tails, and draw a labeled step graph. Use open and filled endpoints correctly and preserve actual monetary spacing. \sscore{\sml{1}}
\end{enumerate}

\columnbreak
\SubsectionBox{Warranty service payouts\score{C2-\sml{2}}}
\QuestionLabel{q:4}
A warranty provider classifies one claim in exactly one of eight exhaustive, mutually exclusive categories: no service, remote diagnosis, software reset, minor repair, major repair, replacement, store credit or cash refund. These eight named categories are the distinct outcomes in the sample space. Their probabilities are 0.20, 0.10, 0.15, 0.15, 0.10, 0.15, 0.05 and 0.10, respectively; do not assume equal likelihood. The corresponding payouts are 0, 2, 4, 6, 10, 15, 15 and 25 USD, respectively. Count store credit at its face value as a payout. Each payout is fixed for its category, and no additional payments apply. Record the provider's payout for one claim.

\begin{enumerate}[label=(\alph*)]
\item Define the random variable and support; derive the grouped PMF; produce its table and a labeled graph with actual monetary spacing. \sscore{\sml{1}}
\item Construct the complete CDF table, including both tails, and draw a labeled step graph. Use open and filled endpoints correctly and preserve actual monetary spacing. \sscore{\sml{1}}
\end{enumerate}

\SubsectionBox{Marketplace account credits\score{C2-\sml{4}}}
\QuestionLabel{q:5}
For one purchase, a marketplace awards exactly one account credit: 0, 2, 5, 9, 14 or 20 USD. These are all possible credits and their categories are mutually exclusive. A 2 USD credit is twice as likely as no credit. A 5 USD credit is as likely as no credit and a 2 USD credit combined. A 9 USD credit is twice as likely as a 5 USD credit. A 14 USD credit is half as likely as a 9 USD credit. A 20 USD credit is as likely as the 2 USD and 14 USD credits combined. There are no other payments or adjustments. Determine the probabilities from these likelihood relationships and record the account credit.\sscore{\sml{1}}

\begin{enumerate}[label=(\alph*)]
\item Define the random variable and support. Derive the six probabilities using normalization and produce the PMF table and labeled graph. \sscore{\sml{2}}
\item Construct the complete CDF table, including both tails, and draw a labeled step graph. Use open and filled endpoints correctly and preserve actual monetary spacing. \sscore{\sml{2}}
\end{enumerate}

\columnbreak
\SubsectionBox{Community market prize tickets\score{C2-\sml{6}}}
\QuestionLabel{q:6}
A community market has nine physically distinct prize tickets: four worth 0 USD each, three worth 5 USD each and two worth 12 USD each. Choose a subset of three tickets uniformly without replacement, so every subset of three physically distinct tickets has the same chance. Ignore selection order. First add the face values. Then add one bonus of 7 USD if the selected subset contains at least one 5 USD ticket and at least one 12 USD ticket; otherwise add no bonus. Finally, if this adjusted amount exceeds 25 USD, pay only 25 USD; otherwise pay the adjusted amount. No other payments or deductions apply. Record the final payout.

\begin{enumerate}[label=(\alph*)]
\item Define the random variable and support. Count the distinct-ticket selections, apply the bonus and then the cap, aggregate equal final payouts, and derive the PMF. Produce its table and labeled graph. \sscore{\sml{3}}
\item Construct the complete CDF table, including both tails, and draw a labeled step graph. Use open and filled endpoints correctly and preserve actual monetary spacing. \sscore{\sml{3}}
\end{enumerate}
\end{multicols*}
\end{document}
